\documentclass[11pt]{article}
\usepackage[margin=1in]{geometry}
\usepackage{amsmath,amssymb,amsthm}
\usepackage{array}
\usepackage{parskip}
\usepackage{booktabs}

\newtheorem{theorem}{Theorem}
\newtheorem{lemma}[theorem]{Lemma}
\newtheorem{corollary}[theorem]{Corollary}
\newtheorem{proposition}[theorem]{Proposition}
\theoremstyle{definition}
\newtheorem{definition}[theorem]{Definition}
\theoremstyle{remark}
\newtheorem{remark}[theorem]{Remark}

\newcommand{\F}{\mathbb{F}}
\newcommand{\dperp}{d^{\perp}}

\title{Disproof of conjecture \texttt{00000008419}}
\author{earthking11}
\date{2026-09-13}

\begin{document}
\maketitle

\section*{Verdict: FALSE}

We disprove conjecture \texttt{00000008419} as stated. The refutation is
unconditional and finite: the binary $[5,2,3]$ code has dual distance
$\dperp = 2$, and at $t = 1$ its two minimum-weight codewords have supports
$\{2,3,5\}$ and $\{1,4,5\}$, which are not a $1$-design. Since the
conjecture's hypothesis $\dperp \ge t+1$ holds ($2 \ge 2$) while its conclusion
fails, the first clause of the conjecture is false. The same code also
falsifies the later clause ``the minimal failure of the converse occurs at
$\dperp = 5$'': the failure occurs already at $\dperp = 2$.

\section{The conjecture and its exact wording}

The statement as filed in \texttt{conjectures/00000008419.md} is a compound,
partly garbled conjecture. The English text reads:

\begin{quote}
\textbf{Definition:} The correspondence between covering radius and design
strength of codes: the supports of dual codewords of weight $w$ form
$t$-designs. \textbf{Conjecture:} For a code with dual distance d-perp at least
$t+1$, all minimal codewords form $t$-$(n,w,\lambda)$ designs; the converse of
the Assmus--Mattson theorem holds at weight $4$, and the minimal failure of the
converse occurs at d-perp $= 5$. (converse of the Assmus--Mattson theorem)
\end{quote}

The Chinese text is equivalent; the bilingual original is reproduced verbatim
in \texttt{README.md}. (It is not typeset here, to avoid shipping a PDF with
missing CJK glyphs: this document deliberately uses no CJK package.)

The statement contains three separable assertions:

\begin{enumerate}
\item[(C1)] \emph{(main clause)} If a binary linear code has dual distance
      $\dperp \ge t+1$, then all minimum-weight codewords of the code form
      $t$-$(n,w,\lambda)$ designs.
\item[(C2)] The converse of the Assmus--Mattson theorem holds at weight $w=4$.
\item[(C3)] The minimal failure of the converse occurs at $\dperp = 5$.
\end{enumerate}

Clause (C1) is false, and it is false already for $t=1$ and $\dperp = 2$. Since
(C3) asserts that no failure occurs below $\dperp = 5$, our witness at
$\dperp = 2$ also falsifies (C3). Because the conjecture is the conjunction of
its clauses, exhibiting one false clause refutes it as filed.

\begin{remark}[on ``minimal codewords'']
Two readings of ``minimal codewords'' are natural: the \emph{minimum-weight}
codewords, and the codewords whose support is minimal under inclusion. Both
readings are falsified by the witness below; see
Section~\ref{sec:robust}. We take the minimum-weight reading in the main
argument, since only then do all blocks have a common size $w$ and the phrase
``$t$-$(n,w,\lambda)$ design'' is well defined.
\end{remark}

\section{Background: the Assmus--Mattson theorem}

The Assmus--Mattson theorem is \emph{not} contradicted by the witness; it is
the conjectured converse that fails. For reference, the theorem (binary case)
states: let $C$ be an $[n,k,d]$ code over $\F_2$ with dual distance $\dperp$,
and let $t$ be an integer with $1 \le t \le \dperp - 1$ such that the number of
weights of $C^{\perp}$ in $\{1,\dots,n-t\}$ is at most $d-t$. Then, for every
$w$, the supports of the codewords of $C$ of weight $w$ form a $t$-design.

The extra hypothesis --- few dual weights in a short interval --- cannot be
dropped. For our witness $C$ we have $d = 3$, $\dperp = 2$, and the nonzero
weights of $C^{\perp}$ are $\{2,3,4\}$. For $t = 1$ the AM condition would
require $\#\{w \in \{1,\dots,4\} : w \text{ a weight of } C^{\perp}\} \le d-t$,
i.e.\ $3 \le 2$, which fails. So the AM theorem simply does not apply to this
code, exactly because its weight-count condition is violated. Clause (C1)
proposes to replace that condition by the much weaker hypothesis
$\dperp \ge t+1$, and it is this strengthening that is false.

\section{The main witness: the binary $[5,2,3]$ code}

\paragraph{Convention.} Coordinates are $1,\dots,5$, and a word is written as a
string $b_1b_2b_3b_4b_5$ whose $i$-th character is the entry at coordinate $i$.

\begin{definition}[The code $C$]
Let
\[
C \;=\; \operatorname{span}_{\F_2}\{\,01101,\; 10011\,\}
\;=\; \{\,00000,\; 01101,\; 10011,\; 11110\,\}.
\]
Its two generators have weight $3$ and their sum $11110$ has weight $4$.
\end{definition}

\begin{theorem}\label{thm:main}
For the code $C$ above:
\begin{enumerate}
\item[(i)] the minimum distance is $d = 3$;
\item[(ii)] the dual distance is $\dperp = 2$ (the word $01100$, of weight $2$,
      lies in $C^{\perp}$, and $C^{\perp}$ has no word of weight $1$);
\item[(iii)] the minimum-weight codewords are exactly $01101$ (support
      $\{2,3,5\}$) and $10011$ (support $\{1,4,5\}$).
\end{enumerate}
Taking $t = 1$, the conjecture's hypothesis $\dperp \ge t+1$ holds, since
$2 \ge 2$. Nevertheless the two supports are \emph{not} a $1$-$(5,3,\lambda)$
design: coordinate $5$ belongs to both blocks while each of the coordinates
$1,2,3,4$ belongs to exactly one.
\end{theorem}

\begin{proof}
Direct computation. The nonzero codewords are $01101$, $10011$ and
$11110$, of weights $3$, $3$ and $4$; hence $d = 3$. The generators and their
sum give the orthogonality conditions for $C^{\perp}$:
\[
y \in C^{\perp} \iff
\begin{cases}
y_2 + y_3 + y_5 = 0,\\
y_1 + y_4 + y_5 = 0.
\end{cases}
\]
The word $01100$ satisfies both (it meets the supports of $01101$, $10011$ and
$11110$ in $2$, $0$ and $2$ positions respectively), so $d^{\perp} \le 2$. No
unit vector $e_i$ is orthogonal to all three nonzero codewords (each coordinate
is non-zero on at least one of them), so $C^{\perp}$ has no word of weight $1$;
hence $\dperp = 2$. The weight-$3$ codewords are exactly the two generators, so
their supports are $\{2,3,5\}$ and $\{1,4,5\}$.

For the design claim observe that the multiset of point multiplicities is
\[
\bigl(\mu_1,\dots,\mu_5\bigr) = (1,1,1,1,2),
\]
so the multiplicity of coordinate $5$ is $2$ while that of coordinates
$1$--$4$ is $1$; a $1$-design requires a constant multiplicity $\lambda$. (In
fact $\lambda$ would have to equal $bw/n = 2\cdot 3/5 = 6/5 \notin \mathbb{Z}$,
which is already impossible.)
\end{proof}

\begin{corollary}
Conjecture \texttt{00000008419} is false: at $(t, \dperp) = (1,2)$ clause (C1)
fails, and the failure occurs at $\dperp = 2 < 5$, so clause (C3) fails too.
\end{corollary}

\begin{remark}[relation to the proposed counterexample]
The refutation proposed in the task statement lists the same code as
$\{\,00000, 01111, 10110, 11001\,\}$ with minimum-weight supports
$\{2,3,5\}$ and $\{1,4,5\}$. That listing reads each binary string with the
\emph{last} character as coordinate $1$; reversing every word
($i \mapsto 6-i$) turns it into the present listing
$\{00000, 01101, 10011, 11110\}$ with the same supports and the same
multiplicities. The two descriptions are the same code up to coordinate
reversal, so the proposal's conclusion is correct. We use the coordinate-first
convention throughout for consistency with the tables and the Lean artifact.
\end{remark}

\section{Robustness: further witnesses and alternative readings}
\label{sec:robust}

The failure does not depend on the minimum weight $3$, on $t = 1$, or on
$\dperp = 2$. Table~\ref{tab:witnesses} records four witnesses and one control;
all quantities are computed by \texttt{reproduce.py} from the generator
matrices over $\F_2$.

\begin{table}[h]
\centering
\small
\begin{tabular}{l l c c c l}
\toprule
& generator rows & $d$ & $\dperp$ & $t$ & design property of the \\
& & & & & minimum-weight codewords \\
\midrule
W1 & $01101,\ 10011$ & $3$ & $2$ & $1$ & \textbf{not} a $1$-design; multiplicities $(1,1,1,1,2)$ \\
W2 & $1110001,\ 1001110$ & $4$ & $2$ & $1$ & \textbf{not} a $1$-design; multiplicities $(2,1,1,1,1,1,1)$ \\
W3 & Hamming $[7,4,3]$ & $3$ & $4$ & $3$ & $1$- and $2$-designs, but \textbf{not} a $3$-design \\
W4 & $110001,\ 101010,\ 011100$ & $3$ & $3$ & $2$ & $1$-design, but \textbf{not} a $2$-design \\
\midrule
W5 & extended Hamming $[8,4,4]$ & $4$ & $4$ & $3$ & \emph{is} a $3$-design (control) \\
\bottomrule
\end{tabular}
\caption{Four witnesses (W1--W4) to the failure of clause (C1), and one control
(W5) on which the design property does hold. Every $t$ in the table satisfies
the conjecture's hypothesis $\dperp \ge t+1$.}
\label{tab:witnesses}
\end{table}

\paragraph{Reading (a): $\dperp \ge t+1$ for the $t$ at which the design is
claimed.} This is the reading used above; W1 (with $t=1$) refutes it.

\paragraph{Reading (b): ``minimal codewords'' = minimal-support codewords.}
For W1 all three nonzero codewords have supports minimal under inclusion: the
two weight-$3$ supports $\{2,3,5\}$, $\{1,4,5\}$ and the weight-$4$ support
$\{1,2,3,4\}$; no one contains another. Grouping by weight, the weight-$3$
minimal codewords are again the two generators and are still not a $1$-design.
For W2 no support contains another either, so all three nonzero words (weights
$4,4,6$) are minimal-support; grouped by weight, the weight-$4$ minimal-support
codewords are $1001110$ and $1110001$, with supports $\{1,4,5,6\}$ and
$\{1,2,3,7\}$, which again fail the $1$-design property. So this reading is
false as well.

\paragraph{Reading (c): restrict to weight $4$.} W2 has minimum distance
$d = 4$, $\dperp = 2$, and its two weight-$4$ codewords have supports
$\{1,4,5,6\}$ and $\{1,2,3,7\}$: coordinate $1$ has multiplicity $2$ and the
others multiplicity $1$, so they are not a $1$-design. Hence even a reading that
confines clause (C1) to weight $4$ is false. The control W5 shows that this is
not a defect of the design test: the extended Hamming $[8,4,4]$ code has
$d = \dperp = 4$ and its $14$ weight-$4$ codewords \emph{do} form a
$3$-$(8,4,1)$ design.

\paragraph{Reading (d): $t \ge 2$.} W4 has $\dperp = 3$, so $t = 2$ is
permitted, and its four minimum-weight codewords form a $1$-design but not a
$2$-design. W3 has $\dperp = 4$ and fails at $t = 3$.

No reading of ``$\dperp$ at least $t+1$'' that keeps the phrase
``all minimal codewords form $t$-$(n,w,\lambda)$ designs'' leaves the statement
true: each of the readings above has an explicit counterexample.

\section{Reproducibility}

\texttt{reproduce.py} (Python~3 standard library only) builds each code from
its generator matrix over $\F_2$ using integer bitsets, computes the minimum
and dual distances, enumerates the minimum-weight codewords and their supports,
and tests the $t$-design property (all $t$-subsets contained in the same number
of blocks). It prints \texttt{PASS} and exits $0$ if and only if every check
holds; in particular it verifies

\begin{quote}\ttfamily
W1: d = 3, d\^{}perp = 2, supports \{2,3,5\}, \{1,4,5\},\\
\hspace*{1.2em} multiplicities (1,1,1,1,2), 1-design = False;\\
W1 violates the AM weight condition (3 > 2 for t = 1),\\
W2: minimum weight 4 and d\^{}perp = 2, 1-design = False,\\
W3: d\^{}perp = 4, 3-design = False,\\
W4: d\^{}perp = 3, 2-design = False,\\
W5: control, 3-design = True.
\end{quote}

\section{Formalisation}

The witness W1 is formalised in the accompanying Lean~4 project under
\texttt{lean4/}. It uses core Lean only (\texttt{import Std}, no Mathlib) and
contains no \texttt{sorry}. Words are \texttt{List Bool} of length $5$; XOR is a
hand-written Boolean XOR so that the kernel reducer can evaluate it (rather than
the opaque \texttt{Nat.xor}); the design predicate counts, for every $t$-subset
of the five coordinates, how many blocks contain it, and tests that all counts
agree. The project proves, among others,

\begin{quote}\ttfamily
theorem minWeight\_eq\_three : minWeight = 3\\
theorem dualDistance\_eq\_two : dualDistance = 2\\
theorem support\_gA : (List.finRange 5).filter (support5 gA) = [1,2,4]\\
theorem support\_gB : (List.finRange 5).filter (support5 gB) = [0,3,4]\\
theorem multiplicity\_five : countBlocks blocks [4] = 2\\
theorem multiplicity\_one : countBlocks blocks [0] = 1\\
theorem isTDesign\_one\_false : isTDesign 1 blocks = false\\
theorem conjecture\_00000008419\_false :\\
\hspace*{2em} dualDistance = 2 \(\wedge\) minWeight = 3 \(\wedge\) isTDesign 1 blocks = false\\
\hspace*{2em} \(\wedge\) countBlocks blocks [4] = 2 \(\wedge\) countBlocks blocks [0] = 1
\end{quote}

\noindent Here \texttt{support5 gA = [1,2,4]} encodes the support
$\{2,3,5\}$ and \texttt{support5 gB = [0,3,4]} encodes $\{1,4,5\}$ (the lists
contain zero-based coordinate indices). \texttt{Check.lean} prints
\texttt{\#print axioms} for every theorem; the audit reports at most
\texttt{propext} and never \texttt{sorryAx}.

\end{document}
